\documentclass[border=15pt]{standalone}
\usepackage[siunitx, RPvoltages]{circuitikz}
\usetikzlibrary{positioning,calc,arrows.meta}

% ============== global style ==============
\ctikzset{
  transistors/arrow pos=end,
  bipoles/length=1.2cm,
  resistors/scale=1.0,
  transistors/scale=1.1,
  tripoles/mos style/arrows,
}
\tikzset{
  pin/.style={circle, fill=black, inner sep=1.5pt, outer sep=0},
  busline/.style={blue!50!black, dashed, line width=0.5mm},
  siglin/.style={red!70!black, dashed, line width=0.5mm},
  nlbl/.style={font=\small},
  dlbl/.style={font=\footnotesize, blue!50!black},
  tlbl/.style={font=\scriptsize\itshape},
}

% ============== column centres (x) ==============
\def\xA{ 2.5}     %% PMOS bias generator
\def\xB{ 8.0}     %% NMOS bias generator
\def\xC{13.5}     %% Tail + input pair
\def\xD{19.5}     %% SE output cascode

% Y-levels (consistent across columns)
\def\yVDD{14.0}                    %% VDD rail
\def\yTopGate{12.7}                %% VP1 bus + top transistor gates
\def\yP1{12.2}    \def\yP1d{10.7}  %% MP1 src / drain
\def\yP2{10.7}    \def\yP2d{ 9.2}  %% MP2 src / drain  (also VP2-tap level)
\def\yMid{ 7.5}                    %% Rp/Rn midpoint
\def\yN2{ 6.0}    \def\yN2s{ 4.5}  %% MN2 drain/src  (VN2 tap)
\def\yN1{ 4.5}    \def\yN1s{ 3.0}  %% MN1 drain/src  (VN1 tap)
\def\yGND{ 2.0}

\begin{document}
\begin{circuitikz}[american, font=\small]

% ===================================================================
% TITLE
% ===================================================================
\node[font=\Large\bfseries] at (\xC, 15.2)
   {Folded-cascode SE OTA --- wide-swing R-offset bias};
\node[font=\normalsize] at (\xC, 14.6)
   {sky130A LVT \;\; $V_{DD}=0.9\,\mathrm{V}$ \;\; $I_{REF}=10\,\mu\mathrm{A}$};

% ===================================================================
% VDD RAILS  (per column, short stubs to keep diagram clean)
% ===================================================================
\foreach \x in {\xA,\xB,\xC,\xD}{
  \draw[thick] (\x-1.4,\yVDD) -- (\x+1.4,\yVDD);
  \node[font=\normalsize] at (\x-1.6, \yVDD) {$V_{DD}$};
}

% ===================================================================
% COLUMN A : PMOS bias generator
%   VDD - MP1 (diode VP1) - MP2 (diode VP2) - Rp - 10uA - GND
% ===================================================================

% --- MP1 (PMOS, source up to VDD, diode-tied) ---
\node[pmos, xscale=-1, anchor=source] (MP1) at (\xA,\yP1) {};
\draw[thick] (MP1.source) -- (\xA,\yVDD);
% gate stub on LEFT (xscale flips, so gate ends up on right of -1 = left)
\draw (MP1.gate) -- ++(-0.6,0) coordinate (gP1);
\node[pin] at (gP1) {};
\node[nlbl, left] at (gP1) {$V_{P1}$};
% diode-tie: gate <-> drain via outside loop on LEFT
\draw (gP1) -- ($(gP1)+(0,-0.7)$)
       -- ($(MP1.drain)+(-0.6,-0.7)$)
       -- ($(MP1.drain)+(-0.6,0)$)
       -- (MP1.drain);
\node[tlbl, right=12pt of MP1.source, yshift=-4pt] {$M_{P1}$};

% --- MP2 (PMOS, source at MP1 drain, diode-tied to VP2) ---
\node[pmos, xscale=-1, anchor=source] (MP2) at (\xA,\yP2) {};
\draw (MP1.drain) -- (MP2.source);
\draw (MP2.gate) -- ++(-0.6,0) coordinate (gP2);
\node[pin] at (gP2) {};
\node[nlbl, left] at (gP2) {$V_{P2}$};
\draw (gP2) -- ($(gP2)+(0,-0.7)$)
       -- ($(MP2.drain)+(-0.6,-0.7)$)
       -- ($(MP2.drain)+(-0.6,0)$)
       -- (MP2.drain);
\node[tlbl, right=12pt of MP2.source, yshift=-4pt] {$M_{P2}$};

% --- Rp + 10 uA sink ---
\draw (MP2.drain) -- (\xA,\yMid+0.4);
\draw (\xA,\yMid+0.4) to[R, l_=$R_p$] (\xA,\yMid-1.6);
\draw (\xA,\yMid-1.6) -- (\xA,\yMid-2.1);
\draw (\xA,\yMid-2.1) to[I, invert] (\xA,\yGND+0.6);
\node[tlbl, right=4pt] at (\xA,\yMid-2.6) {$10\,\mu$A};
\draw (\xA,\yGND+0.6) -- (\xA,\yGND) node[ground]{};

% column caption
\node[font=\bfseries] at (\xA, 1.0) {PMOS bias gen};

% ===================================================================
% COLUMN B : NMOS bias generator
%   VDD - MPN (g=VP1) - Rn - MN2 (diode VN2) - MN1 (diode VN1) - GND
% ===================================================================

% --- MPN (PMOS mirror, gate from VP1 bus) ---
\node[pmos, xscale=-1, anchor=source] (MPN) at (\xB,\yP1) {};
\draw[thick] (MPN.source) -- (\xB,\yVDD);
\draw (MPN.gate) -- ++(-0.6,0) coordinate (gPN);
\node[pin] at (gPN) {};
\node[tlbl, right=12pt of MPN.source, yshift=-4pt] {$M_{PN}$};

% --- Rn between MPN.drain and MN2.drain ---
\draw (MPN.drain) -- (\xB,\yP2-0.2);
\draw (\xB,\yP2-0.2) to[R, l_=$R_n$] (\xB,\yMid-0.5);

% --- MN2 (NMOS, drain up, diode-tied VN2) ---
\node[nmos, anchor=drain] (MN2) at (\xB,\yMid-0.5) {};
\draw (MN2.gate) -- ++(0.6,0) coordinate (gN2);
\node[pin] at (gN2) {};
\node[nlbl, right] at (gN2) {$V_{N2}$};
\draw (gN2) -- ($(gN2)+(0,0.7)$)
       -- ($(MN2.drain)+(0.6,0.7)$)
       -- ($(MN2.drain)+(0.6,0)$)
       -- (MN2.drain);
\node[tlbl, left=14pt of MN2.drain, yshift=-4pt] {$M_{N2}$};

% --- MN1 (NMOS, source down to GND, diode-tied VN1) ---
\node[nmos, anchor=drain] (MN1) at (\xB,\yMid-2.5) {};
\draw (MN2.source) -- (MN1.drain);
\draw (MN1.gate) -- ++(0.6,0) coordinate (gN1);
\node[pin] at (gN1) {};
\node[nlbl, right] at (gN1) {$V_{N1}$};
\draw (gN1) -- ($(gN1)+(0,0.7)$)
       -- ($(MN1.drain)+(0.6,0.7)$)
       -- ($(MN1.drain)+(0.6,0)$)
       -- (MN1.drain);
\node[tlbl, left=14pt of MN1.drain, yshift=-4pt] {$M_{N1}$};

\draw (MN1.source) -- (\xB,\yGND) node[ground]{};

\node[font=\bfseries] at (\xB, 1.0) {NMOS bias gen};

% ===================================================================
% COLUMN C : Tail + PMOS input pair
%   VDD - MTL(g=VP1) - Vtail - MIP/MIN - Vmidn/Vmidp
% ===================================================================

% --- MTL tail PMOS ---
\node[pmos, xscale=-1, anchor=source] (MTL) at (\xC,\yP1) {};
\draw[thick] (MTL.source) -- (\xC,\yVDD);
\draw (MTL.gate) -- ++(-0.6,0) coordinate (gTL);
\node[pin] at (gTL) {};
\node[tlbl, right=12pt of MTL.source, yshift=-4pt] {$M_{TL}$};

% Vtail node
\draw (MTL.drain) -- (\xC,\yP2-0.3);
\node[pin] at (\xC,\yP2-0.3) {};
\node[nlbl, above right=-2pt and 2pt] at (\xC,\yP2-0.3) {$V_{tail}$};

% Horizontal rail to two input-pair sources
\draw (\xC-2.2,\yP2-0.3) -- (\xC+2.2,\yP2-0.3);

% --- MIP (left of pair) ---  PMOS, source up
\node[pmos, xscale=-1, anchor=source] (MIP) at (\xC-2.2,\yP2-0.3) {};
\draw (MIP.gate) -- ++(-0.7,0) coordinate (gIP);
\node[pin] at (gIP) {};
\node[nlbl, left] at (gIP) {$V_{inp}$};
\node[tlbl, right=12pt of MIP.source, yshift=-4pt] {$M_{INP}$};

% --- MIN (right of pair) --- PMOS, source up
\node[pmos, anchor=source] (MIN) at (\xC+2.2,\yP2-0.3) {};
\draw (MIN.gate) -- ++(0.7,0) coordinate (gIN);
\node[pin] at (gIN) {};
\node[nlbl, right] at (gIN) {$V_{inn}$};
\node[tlbl, left=12pt of MIN.source, yshift=-4pt] {$M_{INN}$};

% Drains drop down: Vmidn (left) and Vmidp (right)
\draw (MIP.drain) -- (\xC-2.2,\yMid-0.7);
\node[pin] at (\xC-2.2,\yMid-0.7) {};
\node[nlbl, below] at (\xC-2.2,\yMid-0.9) {$V_{midn}$};

\draw (MIN.drain) -- (\xC+2.2,\yMid-0.7);
\node[pin] at (\xC+2.2,\yMid-0.7) {};
\node[nlbl, below] at (\xC+2.2,\yMid-0.9) {$V_{midp}$};

\node[font=\bfseries] at (\xC, 1.0) {Tail + input pair};

% ===================================================================
% COLUMN D : SE output cascode (one leg shown)
%   VDD - MP0(g=VP1) - MP2o(g=VP2) - VOUTP - MN2o(g=VN2) - MN1o(g=Vmidn) - GND
% ===================================================================

% --- MP0 top current source ---
\node[pmos, xscale=-1, anchor=source] (MP0) at (\xD,\yP1) {};
\draw[thick] (MP0.source) -- (\xD,\yVDD);
\draw (MP0.gate) -- ++(-0.6,0) coordinate (gP0);
\node[pin] at (gP0) {};
\node[tlbl, right=12pt of MP0.source, yshift=-4pt] {$M_{P0}$};

% --- MP2o cascode PMOS ---
\node[pmos, xscale=-1, anchor=source] (MP2o) at (\xD,\yP2) {};
\draw (MP0.drain) -- (MP2o.source);
\draw (MP2o.gate) -- ++(-0.6,0) coordinate (gP2o);
\node[pin] at (gP2o) {};
\node[tlbl, right=12pt of MP2o.source, yshift=-4pt] {$M_{P2,o}$};

% VOUTP node
\draw (MP2o.drain) -- (\xD,\yP2d-0.5);
\node[pin] at (\xD,\yP2d-0.5) {};
\node[nlbl, right=6pt, font=\bfseries] at (\xD,\yP2d-0.5) {$V_{OUTP}$};

% --- MN2o NMOS cascode ---
\node[nmos, anchor=drain] (MN2o) at (\xD,\yP2d-0.5) {};
\draw (MN2o.gate) -- ++(0.6,0) coordinate (gN2o);
\node[pin] at (gN2o) {};
\node[tlbl, left=14pt of MN2o.drain, yshift=-4pt] {$M_{N2,o}$};

% --- MN1o NMOS bottom (gate driven by Vmidn -- signal) ---
\node[nmos, anchor=drain] (MN1o) at (\xD,\yN2s-1.7) {};
\draw (MN2o.source) -- (MN1o.drain);
\draw (MN1o.gate) -- ++(0.6,0) coordinate (gN1o);
\node[pin] at (gN1o) {};
\node[tlbl, left=14pt of MN1o.drain, yshift=-4pt] {$M_{N1,o}$};

\draw (MN1o.source) -- (\xD,\yGND) node[ground]{};

\node[font=\bfseries] at (\xD, 1.0) {SE output cascode};

% ===================================================================
%  BIAS BUSES  (dashed blue, run at clean y-tracks)
% ===================================================================

% VP1 bus runs at y=\yTopGate, from col A to col D, taps to each gate
\draw[busline] (gP1) -- (gP1 |- \xA,\yTopGate)
   -- (gP0 |- \xD,\yTopGate);
\foreach \g in {gPN,gTL,gP0}{
  \draw[busline] (\g |- \xA,\yTopGate) -- (\g);
}
\node[dlbl, above] at ({(\xA+\xB)/2}, \yTopGate+0.05) {$V_{P1}$};
\node[dlbl, above] at ({(\xB+\xC)/2}, \yTopGate+0.05) {$V_{P1}$};
\node[dlbl, above] at ({(\xC+\xD)/2}, \yTopGate+0.05) {$V_{P1}$};

% VP2 bus runs at y=\yP2d-0.9 from col A (gP2 left of MP2) to col D (gP2o)
% jog: from gP2 go LEFT and DOWN to a clean track at y=\yP2d-0.9, traverse, then up to gP2o
\draw[busline] (gP2) -- ($(gP2)+(-0.5,0)$) coordinate (vp2A)
              -- (vp2A |- \xA,\yP2d-1.0) coordinate (vp2A2)
              -- (vp2A2 -| \xD-1.6, 0)   coordinate (vp2D1)
              -- (vp2D1 |- gP2o)         coordinate (vp2D2)
              -- (gP2o);
\node[dlbl, above] at (\xC, \yP2d-0.95) {$V_{P2}$ bus};

% VN2 bus from col B gate (gN2) to col D MN2o gate (gN2o)
% jog through y = \yN2s-0.2
\draw[busline] (gN2) -- ($(gN2)+(0.6,0)$) coordinate (vn2B)
              -- (vn2B |- \xB,\yN2s-0.4) coordinate (vn2B2)
              -- (vn2B2 -| \xD+1.6,0)    coordinate (vn2D1)
              -- (vn2D1 |- gN2o)         coordinate (vn2D2)
              -- (gN2o);
\node[dlbl, above] at (\xC, \yN2s-0.35) {$V_{N2}$ bus};

% Vmidn SIGNAL from col C left drain to col D MN1o gate
\draw[siglin] (\xC-2.2,\yMid-0.7) -- (\xC-2.2,\yMid-1.4)
              -- (\xD+2.2,\yMid-1.4) -- (\xD+2.2,gN1o |- 0,0)
              -- (gN1o);
\node[dlbl, red!70!black, above] at ({(\xC+\xD)/2}, \yMid-1.35) {$V_{midn}$ \;(signal)};

% ===================================================================
% Foot-note
% ===================================================================
\node[draw, fill=yellow!10, rounded corners, align=left, font=\footnotesize,
      inner sep=6pt] at (\xC, -0.8)
   {\textbf{Topology summary} \quad (matches user's hand-drawn diagram)\\[2pt]
    \textbullet\ Column A generates $V_{P1}, V_{P2}$ via two PMOS diodes + $R_p$ + ext.\ 10\,$\mu$A sink.\\
    \textbullet\ Column B generates $V_{N1}, V_{N2}$ via $R_n$ + two NMOS diodes, mirrored from $V_{P1}$.\\
    \textbullet\ Column C: PMOS tail (g=$V_{P1}$) + PMOS input pair, drains $\to V_{midn}, V_{midp}$.\\
    \textbullet\ Column D: SE folded-cascode output; $M_{P0}|M_{P2,o}$ supply, $M_{N2,o}|M_{N1,o}$ sink,\\
    \quad bottom NMOS gate driven by $V_{midn}$ from the input pair (the folding signal).};

\end{circuitikz}
\end{document}
